\subsection{Two double hom-sets force a nontrivial graded graph}
\begin{theorem}[Proved]\label{A7:two-double-graph}
Let $C$ be a finite acyclic category with support $0<1<2<3<4$.
Suppose $C(0,1)=\{a_0,a_1\}$ and $C(1,2)=\{b_0,b_1\}$.
Define a directed multigraph $\Gamma$ with vertex set $C(1,3)$,
and, for each $c\in C(2,3)$, one edge from $c\circ b_1$ to
$c\circ b_0$. If some product of rational degree-two classes on
$BC$ is nonzero, the following conditions are necessary:
\begin{enumerate}
\item There are distinct $u_0,u_1\in C(0,2)$ with
$b_j\circ a_i=u_i$ for every $i,j\in\{0,1\}$.
\item There is a function $h:C(1,3)\to\mathbb Q$ satisfying
$h(c\circ b_0)-h(c\circ b_1)=1$ for every $c\in C(2,3)$.
\item After identifying parallel edges, the undirected graph
underlying $\Gamma$ contains a cycle. Every such cycle has even
length, and hence length at least four.
\end{enumerate}
In particular, if $|C(2,3)|\le3$ or $|C(1,3)|\le3$, all rational
degree-two products vanish. The dual statement holds with the
roles of $012$ and $234$ reversed.
\end{theorem}
\begin{proof}
Write $A=C(0,1)$, $B=C(1,2)$, $D=C(2,3)$ and $E=C(3,4)$, and set
$u(a,b)=b\circ a$, $v(b,c)=c\circ b$, $w(c,d)=d\circ c$.
Let $\alpha,\beta$ be two-cocycles with nonzero product. The
restriction of $[\alpha]$ to the full subcategory on $012$ cannot
be zero: otherwise a primitive one-cochain extended by zero makes
$\alpha$ vanish on all $012$ triangles, annihilating its product
with $\beta$ by the Alexander--Whitney formula.

Every two-cycle on $012$ is a multiple of
$(a_0-a_1)\otimes(b_0-b_1)$. Thus
\[
 [u(a_0,b_0)]+[u(a_1,b_1)]
 =[u(a_0,b_1)]+[u(a_1,b_0)],
\]
and, with $\ell_b=\alpha(a_0,b)-\alpha(a_1,b)$, we have
$x=\ell_{b_0}-\ell_{b_1}\ne0$. Equality of these multisets gives
exactly two disjoint alternatives: either $u(a_0,b)=u(a_1,b)$
for each $b$ (case U0), or $u(a,b_0)=u(a,b_1)$ for each $a$
with distinct row values (case U1).

Define $F:C(1,3)\to\mathbb Q$ by
$F(s)=\alpha(a_0,s)-\alpha(a_1,s)$. The two cocycle equations
on $(a_0,b,c)$ and $(a_1,b,c)$, subtracted, show that
\[
 \ell_b=F(v(b,c))\quad\text{in U0},\qquad
 x=F(v(b_0,c))-F(v(b_1,c))\quad\text{in U1}.
\]
Let $z$ be any rational four-cycle. Its face omitting $0$ forces
\[
 z=(a_0-a_1)\otimes
   \sum_{b\in B,c\in D,d\in E}t_{bcd}\,b\otimes c\otimes d.
\]
The four remaining separate face summands give
\begin{align*}
 &\sum_d t_{bcd}=0 &&(b\in B,c\in D),\\
 &\sum_b t_{bcd}([u(a_0,b)]-[u(a_1,b)])=0 &&(c\in D,d\in E),\\
 &\sum_{b,c}t_{bcd}[v(b,c)]=0 &&(d\in E),\\
 &\sum_{c,d}t_{bcd}[w(c,d)]=0 &&(b\in B).
\end{align*}
All equalities involving brackets are in the free rational vector
space on the indicated hom-set.

In U0, the images $v(b_0,D)$ and $v(b_1,D)$ are disjoint: a common
value would give $\ell_{b_0}=\ell_{b_1}$, contrary to $x\ne0$.
The third boundary equality therefore holds separately for each
$b$. The cocycle equation on $(b,c,d)$ is
\[
 \beta(c,d)=\beta(v(b,c),d)-\beta(b,w(c,d))+\beta(b,c).
\]
Multiplying by $t_{bcd}$ and summing over $c,d$, all three terms
on the right vanish: respectively by the separated third boundary
equality, the fourth equality, and the first equality. Therefore
\[
 \langle\alpha\smile\beta,z\rangle
 =\sum_b\ell_b\sum_{c,d}t_{bcd}\beta(c,d)=0.
\]
Since this holds for every $z$, U0 contradicts the nonzero product.
Consequently U1 holds, proving condition (1).

In U1 the second boundary equality forces
$t_{b_0cd}=q_{cd}$ and $t_{b_1cd}=-q_{cd}$ for scalars $q_{cd}$.
The remaining equalities give
\[
 \sum_dq_{cd}=0,\qquad
 \sum_cq_{cd}([v(b_0,c)]-[v(b_1,c)])=0,\qquad
 \sum_{c,d}q_{cd}[w(c,d)]=0.
\]
The function $h=F/x$ satisfies condition (2). In each connected
component, its values differ from one fixed rational offset by
integers, with each edge increasing the integer by one. Hence the
underlying graph is bipartite; there are no loops and opposite
orientations on the same pair of vertices cannot both occur.

Suppose its underlying simple graph were a forest. The kernel of
its incidence matrix is then generated by differences of parallel
edges. More explicitly, orient each edge by increasing $h$.
The coefficients of a one-cycle sum to zero over each parallel
class: remove a leaf of each finite tree and use the incidence
equation at that leaf, then continue inductively. Applied to the
middle displayed equality, this says that for every $d$ and every
parallel class $P\subseteq D$, $\sum_{c\in P}q_{cd}=0$.

Using the cocycle equation for $b=b_0$ and the first and third
remaining boundary equalities gives
\[
 \sum_{c,d}q_{cd}\beta(c,d)
 =\sum_{c,d}q_{cd}\beta(v(b_0,c),d)=0.
\]
The last equality follows because $v(b_0,c)$ is constant on each
parallel class. But the cup evaluation is $x$ times the left side,
so again it vanishes on every four-cycle, a contradiction. This
proves condition (3). A bipartite simple cycle requires at least
four edges and four vertices, proving the numerical consequences.
Opposition reverses the five objects and proves the dual statement.
\end{proof}
